\begin{center}\begin{tikzpicture}[font=\small,>=stealth,line join=round]
\draw[gray!40,rounded corners=3pt] (0,0) rectangle (14.2,8.5);\node at (7.1,7.9) {易错一：关于原点对称，不等于图像一定包含原点};
\node at (3.5,6.85){$y=x^3$：$0$ 可以代入};\node at (10.6,6.85){$y=1/x$：$0$ 不能代入};\begin{scope}[shift={(3.5,3.7)},x=1cm,y=1cm]\draw[->,gray] (-2.2,0)--(2.2,0) node[right]{$x$};\draw[->,gray] (0,-2.2)--(0,2.2) node[above]{$y$};\draw[blue,thick,domain=-1.25:1.25,samples=150] plot(\x,{(\x)^3});\fill[orange!85!black] (0,0) circle(2.2pt) node[below right=5pt] {$O$};\end{scope}\begin{scope}[shift={(10.5,3.7)},x=1cm,y=1cm]\draw[->,gray] (-2.2,0)--(2.2,0) node[right]{$x$};\draw[->,gray] (0,-2.2)--(0,2.2) node[above]{$y$};\draw[blue,thick,domain=-2:-0.5,samples=150] plot(\x,{1/\x});\draw[blue,thick,domain=0.5:2,samples=150] plot(\x,{1/\x});\fill[orange!85!black] (1,1) circle(2.2pt) node[above right=5pt] {$(1,1)$};\fill[orange!85!black] (-1,-1) circle(2.2pt) node[below left=5pt] {$(-1,-1)$};\node[below right=5pt] at (0,0){$O$};\end{scope}\node at (3.5,1.25){是奇函数，且 $f(0)=0$};\node at (10.6,1.25){是奇函数，但 $f(0)$ 不存在};
\node[text width=12.8cm,align=center] at (7.1,.65) {两图横纵同尺度；右图原点仅是坐标轴交点，不是函数图像上的点。};
\end{tikzpicture}\end{center}
